% 辅助线方法 A：过 A 作 AE∥DC 交 BC 于 E（平移一腰）
\documentclass[tikz,border=4pt]{standalone}
\usepackage{ctex}
\usepackage{amsmath}
\usetikzlibrary{calc}
\begin{document}
\begin{tikzpicture}[scale=1.0, thick]
  \coordinate (A) at (1.2,2.5);
  \coordinate (D) at (3.8,2.5);
  \coordinate (B) at (0,0);
  \coordinate (C) at (5,0);
  % AE ∥ DC，故 E = B + (D-A) 平移到 BC 上：方向向量 (DC)
  % E 在 BC 上：从 A 出发沿 DC 方向 (5-3.8, 0-2.5) = (1.2,-2.5) 到 BC 即 y=0
  % 已知 A=(1.2,2.5), 沿 (1.2,-2.5) 走 t=1 到 (2.4,0)
  \coordinate (E) at (2.4,0);

  \draw (A)--(B)--(C)--(D)--cycle;
  \draw[red] (A)--(E);

  \node[above left]  at (A) {$A$};
  \node[above right] at (D) {$D$};
  \node[below left]  at (B) {$B$};
  \node[below right] at (C) {$C$};
  \node[below] at (E) {$E$};

  \node[red, right] at (1.8,1.3) {\footnotesize $AE \parallel DC$};
  \node[below] at (2.5,-0.6) {\footnotesize 方法 A：平移一腰};
\end{tikzpicture}
\end{document}
